\section{Consistency and gradient-flow proofs}
\label{app:consistency}

\subsection{Proof of \hyperref[lem:expansions]{Lemma 1}}
\label{app:expansions}

\begin{proof}
For a square matrix $\Gamma$ and $h\ge0$, the integral remainder formula for
the matrix exponential gives
\begin{equation}
 \e^{-h\Gamma}
 =I-h\Gamma+
 h^2\int_0^1(1-\theta)\Gamma^2\e^{-\theta h\Gamma}\dd\theta.
 \label{eq:matrix-exp-remainder}
\end{equation}
Consequently,
\begin{equation}
 \norm{\e^{-h\Gamma}-I+h\Gamma}
 \le \frac{h^2}{2}\norm{\Gamma}^2\e^{h\norm{\Gamma}}.
 \label{eq:matrix-exp-bound}
\end{equation}
If $\norm{\Gamma}\le M_\Gamma$ on the coefficient set under consideration,
the right-hand side is bounded by
$\tfrac12M_\Gamma^2\e^{h_0M_\Gamma}h^2$ for $0\le h\le h_0$.

Similarly, for $\eta\ge0$,
\begin{equation}
 1-\e^{-h\eta}
 =h\eta-h^2\eta^2\int_0^1(1-\theta)\e^{-\theta h\eta}\dd\theta,
 \label{eq:beta-remainder}
\end{equation}
so that
\begin{equation}
 \left|1-\e^{-h\eta}-h\eta\right|
 \le\frac12h^2\eta^2.
 \label{eq:beta-remainder-bound}
\end{equation}
Uniform boundedness of $\eta$ gives a uniform $O(h^2)$ remainder. Applying
these estimates with $h=h_n$, $\Gamma=\Gamma_n$, and $\eta=\eta_n$ proves
\eqref{eq:first-order-expansions}.
\end{proof}

\subsection{Proof of \hyperref[thm:local-consistency]{Theorem 1}}
\label{app:local-consistency}

\begin{proof}
All quantities in this proof are evaluated at $t_n$ unless stated otherwise.
Write
\begin{equation}
 D=I-h\Gamma+R_D,\qquad
 \beta=h\eta+r_\beta,
 \label{eq:local-remainders}
\end{equation}
where \hyperref[lem:expansions]{Lemma 1} gives
\begin{equation}
 \norm{R_D}\le c_Dh^2,\qquad
 |r_\beta|\le c_\beta h^2.
 \label{eq:local-remainder-bounds}
\end{equation}
With $\overline S=DS$,
\begin{equation}
 \overline S=S-h\Gamma S+R_DS.
 \label{eq:sbar-expansion}
\end{equation}
The reconstruction error evaluated after forgetting satisfies
\begin{align}
 v-\overline S^\top k
 &=
 v-S^\top k+hS^\top\Gamma^\top k-(R_DS)^\top k \notag\\
 &=v-S^\top k+O(h),
 \label{eq:forgotten-error-expansion}
\end{align}
uniformly on a compact set of $(S,k,v,\Gamma)$. Substitution into
\eqref{eq:erase-write} yields
\begin{align}
 S^{\rm rec}_{n+1}
 &=
 S-h\Gamma S+R_DS
 +(h\eta+r_\beta)k
 \left[v-S^\top k+O(h)\right]^\top \notag\\
 &=
 S+h\left[-\Gamma S+\eta k(v-S^\top k)^\top\right]
 +R_{\rm rec},
 \label{eq:recurrence-local-expansion}
\end{align}
where $\norm{R_{\rm rec}}\le C_{\rm rec}h^2$. The continuous vector field at
$t_n$ is precisely the bracketed term. Hence
\begin{equation}
 \fnorm{
 S^{\rm rec}_{n+1}-S(t_n)-h\dot S(t_n)}
 \le C_{\rm rec}h^2.
 \label{eq:recurrence-defect-proof}
\end{equation}
The assumed solution regularity also gives
\begin{equation}
 \fnorm{S(t_{n+1})-S(t_n)-h\dot S(t_n)}
 \le C_{\rm sol}h^2
 \label{eq:solution-taylor-remainder}
\end{equation}
with a uniform $C_{\rm sol}$. The triangle inequality proves
\eqref{eq:local-defect} with
$C_{\rm loc}=C_{\rm rec}+C_{\rm sol}$.
\end{proof}

\subsection{Proof of \hyperref[cor:global-first-order]{Corollary 1}}
\label{app:global-convergence}

\begin{proof}
Let $\Phi_n$ denote one step of the recurrence and set
$e_n=\fnorm{S_n-S(t_n)}$. Add and subtract
$\Phi_n(S(t_n))$:
\begin{align}
 e_{n+1}
 &\le
 \fnorm{\Phi_n(S_n)-\Phi_n(S(t_n))}
 +\fnorm{\Phi_n(S(t_n))-S(t_{n+1})}.
 \label{eq:error-recursion-start}
\end{align}
The one-step Lipschitz assumption bounds the first term by
$(1+L_hh_n)e_n$.
\hyperref[thm:local-consistency]{Theorem 1} bounds the second term by
$C_{\rm loc}h_n^2$. Since $h_n^2\le hh_n$,
\begin{equation}
 e_{n+1}\le(1+L_hh_n)e_n+C_{\rm loc}hh_n.
 \label{eq:error-recursion}
\end{equation}
For identical initial data, $e_0=0$. Iteration and
$1+x\le\e^x$ give
\begin{align}
 e_n
 &\le C_{\rm loc}h
 \sum_{j=0}^{n-1}h_j
 \prod_{\ell=j+1}^{n-1}(1+L_hh_\ell) \notag\\
 &\le C_{\rm loc}h
 \sum_{j=0}^{n-1}h_j
 \e^{L_h(t_n-t_{j+1})}.
 \label{eq:discrete-gronwall-sum}
\end{align}
For $L_h>0$, the right Riemann sum is bounded by
\begin{equation}
 \frac{C_{\rm loc}}{L_h}
 (\e^{L_ht_n}-1)h
 \le
 \frac{C_{\rm loc}}{L_h}
 (\e^{L_hT}-1)h.
 \label{eq:discrete-gronwall-result}
\end{equation}
If $L_h=0$, \eqref{eq:error-recursion} gives
$e_n\le C_{\rm loc}t_nh\le C_{\rm loc}Th$.
\end{proof}

\subsection{Proof of \hyperref[prop:gradient-flow]{Proposition 1}}
\label{app:gradient-flow}

\begin{proof}
Let $R(S)=v-S^\top k\in\R^{d_v}$. For a variation
$H\in\R^{d_k\times d_v}$,
\begin{align}
 \left.\frac{\dd}{\dd\epsilon}
 \frac12\norm{R(S+\epsilon H)}^2\right|_{\epsilon=0}
 &=-R(S)^\top H^\top k \notag\\
 &=-\tr\!\left([kR(S)^\top]^\top H\right).
 \label{eq:first-objective-variation}
\end{align}
Thus the gradient of the first term is $-kR(S)^\top$. Since $\Gamma$ is
symmetric,
\begin{align}
 \left.\frac{\dd}{\dd\epsilon}
 \frac{1}{2\eta}
 \tr[(S+\epsilon H)^\top\Gamma(S+\epsilon H)]
 \right|_{\epsilon=0}
 =\frac1\eta\tr[(\Gamma S)^\top H].
 \label{eq:regularizer-variation}
\end{align}
Therefore
\begin{equation}
 \nabla_S\mathcal J
 =-k(v-S^\top k)^\top+\frac1\eta\Gamma S.
 \label{eq:objective-gradient}
\end{equation}
Multiplication by $-\eta$ produces
\eqref{eq:continuous-memory}.
\end{proof}

\subsection{Proof of \hyperref[cor:fixed-equilibrium]{Corollary 2}}
\label{app:equilibrium}

\begin{proof}
With fixed coefficients, \eqref{eq:continuous-memory} is
\begin{equation}
 \dot S=-(\Gamma+\eta kk^\top)S+\eta kv^\top.
 \label{eq:fixed-memory-linear}
\end{equation}
Set $A=\Gamma+\eta kk^\top$. Since $\Gamma\succ0$ and $\eta kk^\top$
is positive semidefinite, $A$ is symmetric positive definite. Therefore
$A^{-1}$ exists, and the unique equilibrium is
$S_\star=A^{-1}\eta kv^\top$.

The error $E=S-S_\star$ obeys $\dot E=-AE$, hence
$E(t)=\e^{-At}E(0)$. Symmetry and positive definiteness imply
\begin{equation}
 \norm{\e^{-At}}\le\e^{-\lambda_{\min}(A)t}.
 \label{eq:spd-semigroup}
\end{equation}
The Frobenius norm is subordinate column by column to the spectral norm,
which proves \eqref{eq:equilibrium-convergence}.
\end{proof}

\section{Boundedness and well-posedness proofs}
\label{app:boundedness}

\subsection{Proof of \hyperref[thm:local-wellposedness]{Theorem 2}}
\label{app:wellposedness}

\begin{proof}
Identify $S\in\R^{d_k\times d_v}$ with
$s=\operatorname{vec}(S)\in\R^{d_kd_v}$. On a bounded set in $(a,S)$,
A2 bounds the local Lipschitz constant of $F_G$. By A3, the coefficient maps
$k,v,q,\Gamma,\eta$ are locally Lipschitz and bounded on that set.

The map $(a,S)\mapsto S^\top q(a,\mu)$ is locally Lipschitz because, for two
states,
\begin{align}
 \norm{S_1^\top q(a_1)-S_2^\top q(a_2)}
 &\le
 \norm{S_1-S_2}\norm{q(a_1)}
 +\norm{S_2}\norm{q(a_1)-q(a_2)}.
 \label{eq:read-lipschitz}
\end{align}
The same product estimate applies to
$\Gamma(a,\mu)S$ and to
$\eta(a,\mu)k(a,\mu)[v(a,\mu)-S^\top k(a,\mu)]^\top$.
Consequently, the vectorized right-hand side of \eqref{eq:cdm-grom} is
locally Lipschitz in $(a,s)$. Picard--Lindel\"of yields a unique solution on a
nonzero interval and a unique maximal continuation.

No linear-growth or dissipativity condition on the full augmented vector
field has been assumed. The maximal interval therefore need not be infinite,
which is why the statement is local.
\end{proof}

\subsection{Proof of \hyperref[thm:memory-boundedness]{Theorem 3}}
\label{app:memory-bound}

\begin{proof}
Set $y=S^\top k\in\R^{d_v}$. Taking the Frobenius product of
\eqref{eq:continuous-memory} with $S$ gives
\begin{align}
 \frac12\frac{\dd}{\dd t}\fnorm S^2
 &=-\ipF{S}{\Gamma S}
 +\eta\ipF{S}{k(v-S^\top k)^\top} \notag\\
 &=-\ipF{S}{\Gamma S}
 +\eta y^\top(v-y).
 \label{eq:memory-energy-full}
\end{align}
The last equality follows from
\begin{equation}
 \ipF{S}{k(v-y)^\top}
 =\tr(S^\top k(v-y)^\top)
 =y^\top(v-y).
 \label{eq:frobenius-write}
\end{equation}
Complete the square:
\begin{equation}
 y^\top(v-y)
 =-\norm{y-\tfrac12v}^2+\frac14\norm v^2
 \le\frac14\norm v^2.
 \label{eq:complete-square}
\end{equation}
Using A5--A7 in \eqref{eq:memory-energy-full},
\begin{equation}
 \frac12\frac{\dd}{\dd t}\fnorm S^2
 \le
 -\gamma_{\min}\fnorm S^2+
 \frac{\eta_{\max}v_{\max}^2}{4}.
 \label{eq:scalar-memory-inequality}
\end{equation}
Let $z(t)=\fnorm{S(t)}^2$. Then
\begin{equation}
 z'(t)+2\gamma_{\min}z(t)
 \le\frac{\eta_{\max}v_{\max}^2}{2}.
 \label{eq:z-inequality}
\end{equation}
Multiplication by $\e^{2\gamma_{\min}t}$ and integration from $0$ to $t$
give
\begin{equation}
 z(t)\le
 \e^{-2\gamma_{\min}t}z(0)
 +\frac{\eta_{\max}v_{\max}^2}{4\gamma_{\min}}
 (1-\e^{-2\gamma_{\min}t}),
 \label{eq:z-bound-proof}
\end{equation}
which is \eqref{eq:memory-explicit-bound} on the maximal interval of the
coupled solution. If that interval is $[0,\infty)$, taking $\limsup$ proves
\eqref{eq:memory-limsup}.
\end{proof}

\subsection{Proof of \hyperref[prop:memory-passivity]{Proposition 2}}
\label{app:passivity}

\begin{proof}
Identity \eqref{eq:memory-energy-identity} was established in
\eqref{eq:memory-energy-full}. Young's inequality gives
\begin{equation}
 y^\top v
 \le\frac12\norm y^2+\frac12\norm v^2.
 \label{eq:young-yv}
\end{equation}
Therefore
\begin{equation}
 y^\top(v-y)
 \le-\frac12\norm y^2+\frac12\norm v^2.
 \label{eq:passivity-square}
\end{equation}
Use A5 in the forgetting term and multiply
\eqref{eq:passivity-square} by the nonnegative $\eta$ from A6 to obtain
\eqref{eq:passivity-bound}.
\end{proof}

\section{Finite-state, Volterra, and linear-elimination proofs}
\label{app:finite-state}

\subsection{Proof of \hyperref[prop:fixed-key]{Proposition 3}}
\label{app:fixed-key}

\begin{proof}
Under the stated specialization, the matrix equation is
\begin{equation}
 \dot S=-\gamma S+\eta e(Ca-S^\top e)^\top.
 \label{eq:fixed-key-matrix}
\end{equation}
Because $y=S^\top e$ and $\norm e=1$,
\begin{align}
 \dot y
 &=\dot S^\top e \notag\\
 &=-\gamma S^\top e+
 \eta(Ca-S^\top e)e^\top e \notag\\
 &=-(\gamma+\eta)y+\eta Ca,
 \label{eq:fixed-key-proof}
\end{align}
which proves \eqref{eq:fixed-key-vector}.

Let $P_\perp=I-ee^\top$. Since $P_\perp e=0$,
\begin{equation}
 \frac{\dd}{\dd t}(P_\perp S)
 =-\gamma P_\perp S.
 \label{eq:orthogonal-equation}
\end{equation}
Solving this linear equation gives \eqref{eq:orthogonal-decay}. Moreover,
\begin{equation}
 S^\top e=(ee^\top S)^\top e,
 \label{eq:read-projection}
\end{equation}
so $P_\perp S$ is unobservable from the read and closure. The observable
single-address dynamics are therefore represented by $y$. Further reduction
can occur only if the ordinary linear input/output realization defined by
$B$ and $C$ is itself nonminimal.
\end{proof}

\subsection{Proof of \hyperref[thm:volterra]{Theorem 4}}
\label{app:volterra}

\begin{proof}
Variation of constants applied to \eqref{eq:auxiliary-system} gives
\begin{equation}
 y_j(t)=\e^{-\Lambda_jt}y_j(0)
 +\int_0^t\e^{-\Lambda_j(t-s)}C_ja(s)\dd s.
 \label{eq:auxiliary-voc}
\end{equation}
Substitute \eqref{eq:auxiliary-voc} into
$\dot a=F_G(a;\mu)+\sum_jB_jy_j$:
\begin{align}
 \dot a(t)
 &=F_G(a(t);\mu)
 +\sum_jB_j\e^{-\Lambda_jt}y_j(0) \notag\\
 &\quad+
 \int_0^t\sum_j
 B_j\e^{-\Lambda_j(t-s)}C_ja(s)\dd s.
 \label{eq:volterra-substitution}
\end{align}
The finite sum may be interchanged with the integral. Definitions
\eqref{eq:initial-slip} and \eqref{eq:finite-kernel} then give
\eqref{eq:volterra-resolved}.

Exponential stability of $-\Lambda_j$ makes
$\e^{-\Lambda_jt}$ exponentially integrable. For $\operatorname{Re}s$
larger than the semigroup growth bound,
\begin{equation}
 \int_0^\infty\e^{-st}\e^{-\Lambda_jt}\dd t
 =(sI+\Lambda_j)^{-1}.
 \label{eq:matrix-laplace}
\end{equation}
Taking the Laplace transform term by term proves
\eqref{eq:kernel-laplace}.
\end{proof}

\subsection{Proof of \hyperref[thm:linear-elimination]{Theorem 5}}
\label{app:linear-elimination}

\begin{proof}
The unresolved block of \eqref{eq:block-linear} is
\begin{equation}
 \dot c=A_{uu}c+A_{ur}a.
 \label{eq:unresolved-block}
\end{equation}
Variation of constants yields
\begin{equation}
 c(t)=\e^{A_{uu}t}c(0)+
 \int_0^t\e^{A_{uu}(t-s)}A_{ur}a(s)\dd s.
 \label{eq:unresolved-voc}
\end{equation}
The resolved block is
$\dot a=A_{rr}a+A_{ru}c$. Substitution of
\eqref{eq:unresolved-voc} gives
\begin{align}
 \dot a(t)
 &=A_{rr}a(t)+A_{ru}\e^{A_{uu}t}c(0) \notag\\
 &\quad+
 \int_0^tA_{ru}\e^{A_{uu}(t-s)}A_{ur}a(s)\dd s,
 \label{eq:linear-substitution}
\end{align}
which proves \eqref{eq:exact-linear-memory} and
\eqref{eq:exact-kernel}.

This algebra is valid on finite intervals for every finite matrix $A_{uu}$.
Under A8, standard finite-dimensional semigroup theory provides constants
$M_u\ge1$ and $\omega_u>0$ such that
\begin{equation}
 \norm{\e^{A_{uu}t}}\le M_u\e^{-\omega_ut},
 \label{eq:hurwitz-semigroup}
\end{equation}
so the slip and kernel decay exponentially.
\end{proof}

\subsection{Proof of \hyperref[cor:exact-realization]{Corollary 3}}
\label{app:exact-realization}

\begin{proof}
With $\Lambda=-A_{uu}$, $B=A_{ru}$, and $C=A_{ur}$,
\begin{equation}
 B\e^{-\Lambda\tau}C
 =A_{ru}\e^{A_{uu}\tau}A_{ur}
 =K_{\rm exact}(\tau).
 \label{eq:realization-identity}
\end{equation}
The auxiliary initial-slip term is
$B\e^{-\Lambda t}y(0)=A_{ru}\e^{A_{uu}t}c(0)$ when $y(0)=c(0)$.
A8 makes $-\Lambda=A_{uu}$ Hurwitz and hence makes the realization
exponentially stable.

If stable model reduction replaces $(A_{uu},A_{ur},A_{ru})$ by a lower-order
realization, the identity becomes an approximation
$K_{\rm exact}\approx K_m$. The resulting trajectory error is then bounded
under the hypotheses of \hyperref[thm:kernel-error]{Theorem 7}.
\end{proof}

\section{Energy and kernel-error proofs}
\label{app:energy-kernel}

\subsection{Proof of \hyperref[thm:reciprocal-energy]{Theorem 6}}
\label{app:reciprocal-energy}

\begin{proof}
Differentiate \eqref{eq:augmented-energy} along
\eqref{eq:reciprocal-model-assumption}:
\begin{align}
 \dot{\mathcal E}
 &=a^\top F_G(a)
 -\sum_j a^\top C_j^\top y_j
 +\sum_j\frac1{\eta_j}
 y_j^\top(\eta_jC_ja-\lambda_jy_j) \notag\\
 &=a^\top F_G(a)
 -\sum_j a^\top C_j^\top y_j
 +\sum_j y_j^\top C_ja
 -\sum_j\frac{\lambda_j}{\eta_j}\norm{y_j}^2.
 \label{eq:energy-before-cancel}
\end{align}
Since the scalar $a^\top C_j^\top y_j$ equals
$y_j^\top C_ja$, the cross terms cancel exactly. This proves
\eqref{eq:reciprocal-energy-identity}.

If $a^\top F_G(a)\le0$, every remaining memory term is nonpositive, so
$\mathcal E$ is non-increasing. Under \eqref{eq:galerkin-absorbing},
\begin{equation}
 \dot{\mathcal E}
 \le c_0-c_1\norm a^2
 -\sum_j\frac{\lambda_j}{\eta_j}\norm{y_j}^2.
 \label{eq:energy-absorbing-step}
\end{equation}
By the definition \eqref{eq:kappa},
\begin{align}
 c_1\norm a^2
 +\sum_j\frac{\lambda_j}{\eta_j}\norm{y_j}^2
 &\ge
 \kappa\left[
 \frac12\norm a^2+
 \sum_j\frac{1}{2\eta_j}\norm{y_j}^2
 \right] \notag\\
 &=\kappa\mathcal E.
 \label{eq:kappa-coercivity}
\end{align}
Thus $\dot{\mathcal E}\le c_0-\kappa\mathcal E$. The integrating-factor
formula gives \eqref{eq:augmented-absorbing-bound}. Under A2 the reciprocal
augmented vector field is locally Lipschitz. Bound
\eqref{eq:augmented-absorbing-bound} confines every finite-time trajectory
to a bounded set, so the standard continuation criterion for ordinary
differential equations excludes a finite maximal endpoint. The solution is
therefore global.
\end{proof}

\subsection{Proof of \hyperref[thm:kernel-error]{Theorem 7}}
\label{app:kernel-error-proof}

\begin{proof}
Let $e(t)=a(t)-a_m(t)$. Subtract the integral forms of
\eqref{eq:exact-volterra} and \eqref{eq:approx-volterra}, and decompose the
convolution difference as
\begin{align}
 K*a-K_m*a_m
 =K*e+(K-K_m)*a_m.
 \label{eq:convolution-decomposition}
\end{align}
For every $t\in[0,T]$, A2 and the common-ball hypothesis give
\begin{align}
 \norm{e(t)}
 &\le\norm{e(0)}
 +L\int_0^t\norm{e(s)}\dd s
 +\int_0^t\norm{g(s)-g_m(s)}\dd s \notag\\
 &\quad+
 \int_0^t\int_0^s
 \norm{K(s-u)}\norm{e(u)}\dd u\dd s \notag\\
 &\quad+
 \int_0^t\int_0^s
 \norm{K(s-u)-K_m(s-u)}
 \norm{a_m(u)}\dd u\dd s.
 \label{eq:kernel-error-pre-fubini}
\end{align}
All integrands are nonnegative and integrable, so Tonelli's theorem applies.
For the term containing $e$,
\begin{align}
 &\int_0^t\int_0^s
 \norm{K(s-u)}\norm{e(u)}\dd u\dd s \notag\\
 &\quad=
 \int_0^t
 \left[\int_0^{t-u}\norm{K(\tau)}\dd\tau\right]
 \norm{e(u)}\dd u \notag\\
 &\quad\le
 \kappa_T\int_0^t\norm{e(u)}\dd u.
 \label{eq:exact-kernel-error-term}
\end{align}
For the kernel mismatch, use $\norm{a_m(u)}\le M$:
\begin{align}
 &\int_0^t\int_0^s
 \norm{K(s-u)-K_m(s-u)}
 \norm{a_m(u)}\dd u\dd s \notag\\
 &\quad\le
 M\int_0^t(t-\tau)\norm{K(\tau)-K_m(\tau)}\dd\tau \notag\\
 &\quad\le MT\delta_K(T).
 \label{eq:kernel-mismatch-term}
\end{align}
The slip mismatch is bounded by $\delta_g(T)$. Therefore
\begin{equation}
 \norm{e(t)}
 \le A_T+(L+\kappa_T)
 \int_0^t\norm{e(u)}\dd u,
 \label{eq:gronwall-kernel}
\end{equation}
where
\begin{equation}
 A_T=\norm{e(0)}+\delta_g(T)+MT\delta_K(T).
 \label{eq:AT}
\end{equation}
Gronwall's inequality gives
$\norm{e(t)}\le A_T\e^{(L+\kappa_T)t}$. Taking the supremum over
$t\in[0,T]$ proves \eqref{eq:kernel-error-bound}. Finally,
\begin{equation}
 \norm{e(0)}+\delta_g(T)+MT\delta_K(T)
 \le\max\{1,MT\}
 [\norm{e(0)}+\delta_g(T)+\delta_K(T)],
 \label{eq:CT-grouping}
\end{equation}
which yields the stated $C_T$.
\end{proof}
